Football analysts describe players with dozens of per-match statistics and then sort them
into role archetypes such as the ball-playing centre back or the deep-lying playmaker. Whether
such types exist has never been tested, because the evidence normally offered for them is
produced just as readily by data containing no groups at all: a clustering algorithm returns
groups whatever it is given, so a selection criterion, a resampling check and an interpretable
average profile are satisfied by a smooth continuum. We test existence directly on \ArcRows{}
player-seasons from \ArcNSeasons{} seasons of the big five European leagues, described by
\ArcFeatures{} statistics, by calibrating cluster quality against simulated data matched to
the observed sample in size, dimension and correlation but built to contain no clusters. One
division survives, and it is a coarse one. Players separate into a defensive and an attacking
mode of involvement, and that separation exceeds every simulation under all three clusterless
references in every scope, by at least \StrMinZOverall{} standard deviations against the
hardest. Nothing finer survives: beyond two modes the observed separation is
indistinguishable from what structureless data produces. The usual stability evidence cannot
make this distinction, since clusterless data resample at an adjusted Rand index of at least
\StrBootNullMinSingle{}. The two modes are not the listed positions renamed, agreeing with
them at \ArcARIvsPosition{}, with defenders and forwards at opposite poles and midfielders
divided almost evenly between the two. The division reproduces in every season and every
league, and a club's composition along it predicts league points beyond possession out of
sample. Football role space has two poles and a populated middle, and finer taxonomies are
artefacts of the method used to find them.
